The locked evaluation contains 1,500 cases and 54,080 candidate intervals. Family-macro AP is 0.438 for the ratio, 0.307 for PCA and 0.423 for same-flow NLL. The paired differences are +0.131 [+0.093, +0.147] against PCA and +0.015 [+0.005, +0.027] against NLL. Both simultaneous primary intervals are positive, supporting the directional H1 for this defined protocol. The NLL contrast is smaller than the chosen 0.02 practical margin.

This aggregate does not imply improvement in every environment. PCA is stronger on both 32-channel synthetic settings, with negative paired intervals for the ratio. The ratio improves on PCA in S64N and the two real families. SKAB contributes most of the macro difference. Intel remains weak in absolute terms despite its positive comparison. Only 65.0\% of Intel target blocks are eligible, and 37 changed targets remain unscorable misses.

The matched PPCA comparison is +0.054 [-0.007, +0.082]. Its interval spans zero, so H2 is inconclusive. The mixture-PPCA interval also spans zero. The conditional Gaussian control is particularly strong in the linear configuration. Supervised trees exceed the ratio in every configuration. Their macro advantage is 0.140, and the paired ratio-minus-tree interval is [-0.169, -0.114]. Thus the results do not establish the ratio as the strongest available detector when labeled injection mechanisms can be exploited. The author GANF and rerun legacy methods provide additional comparisons in Table~\ref{tab:ap}. Top-one ratio attribution is 0.660 on S32N, 0.065 on Intel and 0.631 on SKAB. Held-out mechanisms substantially reduce AP (Table~\ref{tab:strata}).
